\section{Synthesis and Research Agenda}\label{sec:agenda}
\subsection{Which conclusions depend on which sources}\label{sec:sensitivity}
Source withdrawal follows 33 statements across seven comparisons and 20 removal scenarios (Supplement S2.4). With only published texts, WRA and CHeaT retain query recovery from tool destinations, the decoy/visibility utility mismatch behind Finding~2, and a defined prevention outcome. Finding~1's model-view half, live OHTTP refitting and each transfer axis depend on preprints. The controlled studies add implementation-specific measurements, including cover removal using endpoint timing; they do not replace missing source coverage.

\subsection{Lessons for defense designers and evaluators}
\begin{enumerate}
\item \textbf{Name every retained view.} A tool-side or single-step defense leaves model-link, recipient and later-request views unless it is shown to cover them. Report an adapted attack on each permitted view, including single-link rules available to a joint observer.
\item \textbf{Report privacy and utility from the same intervention.} None of the four mapped query/workflow privacy settings documents both an adapted attacker and matched task utility. Utility must include whole-state correctness and resources, not only task completion.
\item \textbf{Treat timing as part of the cover, and budget for it.} Cover inserted outside the model loop is identifiable by its missing decision time, even when the attacker has no labels and the timestamps carry 200\,ms of jitter. Evaluate every cover against a removal attack. Dispatch it on the planner's own gap and duration distribution, and report the added latency next to bytes.
\item \textbf{Calibrate against simple baselines.} Concurrency, position rules and constant-output detectors can rival learned attacks and costly shaping.
\end{enumerate}

\subsection{Tests that would change a protection decision}
Table~\ref{tab:open-questions} turns the synthesis into four discriminating tests: transfer across execution changes, separation of forwarding effects, cover under business constraints, and inference after event discovery. For the ticket application, progress means meeting both boundaries at once: limit the observations that carry task information, and keep decoys within specified state and resource budgets.

\section{Validity, Ethics and Availability}\label{sec:validity}
\paragraph{Evidence scope}
Index failures and missed query terms limit coverage. Extraction is author-performed with AI assistance and has no independent human reliability estimate. Section~\ref{sec:scope} states which fields the findings depend on.

The cover case uses synthetic actions, local services, reused phrasings and supplied call windows. Its timing analyses reuse the collected executions. Each was frozen before its test evaluation, but the size-view results were already known, so they are retrospective new-observer analyses rather than a fresh confirmation cohort. Timestamps are client receipts on loopback; jitter injection approximates, but does not replace, record-level network capture. Timing-matched dispatch is a receipt-level counterfactual that relies on hidden cover not reaching the planner; a live cohort would also measure queueing effects. Intervals condition on fitted attacks and fixtures. Adverse writes do not estimate production prevalence or WRA's harms.

The forwarding pilot uses category-bearing records and twenty nonempty test clusters. Projection changes content, length and execution order together, and its lessons are descriptive. The hosted retest (Supplement S1) changes backend and network together and uses an unpinned alias. It supports no main-text claim.

\paragraph{Ethics and assistance}
All inputs are synthetic, business actions update isolated fixtures, and capture filters restrict collection to the experiment. No human participants, customer operations or provider vulnerabilities are involved. The hosted retest used the account holder's authorization; provider-specific research terms remain unverified. AI tools assisted extraction, implementation and writing; saved observations and executable analyses support every measurement. Traffic inference can enable profiling as well as monitoring. We therefore release payload-free metadata, features, predictions and offline analyses, and exclude credentials, transcripts, raw captures and third-party full texts.

\paragraph{Availability}
The review package links the 44-setting matrix, transfer relations, eligibility ledger and source locators to each claim. It also includes synthetic planner records and fixtures for state replay, predictions, uncertainty scripts and manifests. Raw-to-metadata reconstruction needs the excluded captures.

\section{Conclusion}
Website-fingerprinting evaluation does not transfer unchanged to LLM Agents. Tool feedback carries a task across independently controlled links, so protection must be judged on every view the attacker retains (Finding~1). Cover executes as application work, so privacy and task utility must come from the same intervention, and no mapped privacy setting yet reports both (Finding~2). Executed cover also inherits the timing of execution. Decoys issued outside the model loop lack the planner's decision latency and can be stripped without labels, while timing-matched dispatch restores the size-only protection at a measurable latency cost (Finding~3). The evaluation procedure makes these costs explicit. A defense should be credited only when residual inference, task and state integrity, bytes and latency are reported on the same executions.
